\section{Контрольные вопросы}

\begin{enumerate}
	\item Что такое фазовое пространство?
	\item Каким уравнением описывается поведение лампового генератора? Как его получить?
	\item Какие переменные подаются на входы осциллографа для исследования фазовой плоскости и почему?
	\item При каких условиях для изучения динамики системы применим метод Ван-дер-Поля? Основные положения метода.
	\item Каков вид укороченных уравнений в комплексных амплитудах и для действительных амплитуды и фазы?
	\item Чем определяется выбор аппроксимации характеристики лампы?
	\item Как исследуется устойчивость состояния равновесия исходных и укороченных уравнений?
	\item Какие возможны режимы возбуждения генератора и как экспериментально их менять?
	\item Что такое бифуркационная диаграмма? Как выглядят диаграммы для различных режимов работы генератора?
	\item Какими будут фазовые портреты на различных участках бифуркационной диаграммы?
	\item Как влияют начальные условия на движение точки по фазовой плоскости?
	\item Как изменяется фазовый портрет при внесении затухания в колебательный контур?
\end{enumerate}
